\subsection{Operating curves, pixel selectivity and edit categories}
\label{sec:v19-selectivity}
Each operating point averages the same 42 sources with a nonempty background.
Eight sources have a full edit mask and are excluded from all coordinates,
including directional CLIP. The 18 configurations give 756 source--output
pairs with positive background MAE and positive total change. Table~\ref{tab:v19-selectivity-points}
uses the five observed settings of each editor and three controls. The segments
in the operating plot follow knob order. Matched-response comparisons instead
use each source's own interval of directional CLIP overlap.

For source $i$, let $e_i$ and $b_i$ denote the mean absolute RGB changes per
edit and background pixel, and let $a_i$ be the fraction of pixels in the edit
mask. Whole-image MAE is $T_i=a_i e_i+(1-a_i)b_i$. Selectivity is
$s_i=e_i/b_i$ and the in-mask share of total change is
$\rho_i=a_i s_i/(1-a_i+a_i s_i)$. Each source contributes once to
$\bar s$ and $\bar\rho$. At scale 1.25, spatial SAM gives
$\bar s=4.8752$, $\bar e/\bar b=3.1998$ and median $s_i=3.8419$.
The mean source share is 0.5149 and the share of pooled change is 0.5219.
Table~\ref{tab:v19-selectivity-points} keeps the aggregation order explicit.

\begin{table*}[t]
\centering\small
\caption{All coordinates, including directional CLIP and selectivity, use the same 42 sources with background pixels. Main Table~1 reports directional CLIP on all 50. Background LPIPS uses either spatial feature-map averaging (FM) or masked inputs (MI). The pixel columns report mean whole-image MAE, the mean within-source ratio, the ratio of regional means and the mean within-source share. A star marks the native setting on each five-point curve.}
\label{tab:v19-selectivity-points}
\begin{tabular}{@{}lrrrrrrrr@{}}
\toprule
Configuration & Setting & LPIPS (FM) & LPIPS (MI) & Dir.\ CLIP & $\bar T$ & $\bar s$ & $\bar e/\bar b$ & $\bar\rho$\\
\midrule
Spatial & 0.25 & 0.0140 & 0.0091 & 0.0026 & 0.0132 & 2.0706 & 1.1961 & 0.3519 \\
Spatial & 0.50 & 0.0163 & 0.0107 & 0.0213 & 0.0149 & 2.4762 & 1.4480 & 0.3824 \\
Spatial & 0.75 & 0.0226 & 0.0148 & 0.0451 & 0.0187 & 3.1110 & 1.8887 & 0.4259 \\
Spatial & 1.00$^{*}$ & 0.0462 & 0.0297 & 0.0918 & 0.0283 & 4.0006 & 2.7419 & 0.4798 \\
Spatial & 1.25 & 0.1033 & 0.0593 & 0.1239 & 0.0458 & 4.8752 & 3.1998 & 0.5149 \\
\midrule
Uniform & 0.25 & 0.0156 & 0.0101 & 0.0047 & 0.0139 & 1.8813 & 1.1539 & 0.3436 \\
Uniform & 0.50 & 0.0239 & 0.0154 & 0.0227 & 0.0183 & 2.0542 & 1.2549 & 0.3559 \\
Uniform & 0.75 & 0.0498 & 0.0312 & 0.0457 & 0.0299 & 2.0018 & 1.3554 & 0.3644 \\
Uniform & 1.00$^{*}$ & 0.1656 & 0.0947 & 0.0950 & 0.0600 & 2.1079 & 1.5433 & 0.3795 \\
Uniform & 1.25 & 0.3268 & 0.2063 & 0.1235 & 0.1094 & 1.8753 & 1.4885 & 0.3782 \\
\midrule
FlowEdit & 12 & 0.0541 & 0.0340 & 0.0069 & 0.0222 & 2.1751 & 1.2694 & 0.3548 \\
FlowEdit & 16 & 0.0751 & 0.0478 & 0.0306 & 0.0284 & 2.6358 & 1.3289 & 0.3635 \\
FlowEdit & 20 & 0.1024 & 0.0664 & 0.0555 & 0.0365 & 3.0320 & 1.4193 & 0.3692 \\
FlowEdit & 24$^{*}$ & 0.1825 & 0.1185 & 0.0878 & 0.0602 & 2.5096 & 1.4755 & 0.3717 \\
FlowEdit & 28 & 0.5450 & 0.3546 & 0.1517 & 0.1996 & 1.8067 & 1.1493 & 0.3633 \\
\midrule
Two constants & --- & 0.0646 & 0.0416 & 0.0517 & 0.0344 & 2.0760 & 1.5423 & 0.3706 \\
Uniform @ S & --- & 0.0634 & 0.0423 & 0.0553 & 0.0343 & 2.0924 & 1.5379 & 0.3700 \\
Spatial @ U & --- & 0.1002 & 0.0552 & 0.1198 & 0.0460 & 4.9660 & 3.0792 & 0.5201 \\
\bottomrule
\end{tabular}
\end{table*}

From scale 0.25 to 1.25 the mean selectivity increase is 2.805 [1.717, 4.130]
for spatial SAM and $-0.006$ [$-0.477$, 0.396] for uniform SAM. The paired
difference between these increases is 2.811 [1.925, 3.917]. At scale 1.25
spatial SAM increases selectivity by 3.000 [2.114, 4.088] and the mean
in-mask share by 0.137 [0.099, 0.175] relative to uniform. Selectivity rises
on 37 of 42 sources. Spatial minus uniform directional CLIP is 0.0004
[$-0.0204$, 0.0202] at this setting, with 21 positive and 21 negative
differences. Background feature-map LPIPS is lower for spatial on all 42
sources, by 0.2235 [0.1851, 0.2634] on average.

At the default setting, two constants reduce whole-image MAE relative to
uniform by 0.0255 [0.0213, 0.0304]. Their mean selectivity difference is
$-0.032$ [$-0.147$, 0.083]. Uniform replay of the spatial profile gives
a mean selectivity of 2.092, close to the two-constant value of 2.076.
Spatial replay of the uniform profile gives 4.966, 2.858 [2.055, 3.846]
above native uniform. These intervals use 10,000 paired source bootstrap
draws with seed 0 and percentile endpoints. The contrasts measure changes
in pixels and text alignment under the specified settings.

\begin{table*}[t]
\centering\small
\caption{Matched uniform minus spatial background LPIPS by PIE-Bench edit category. The denominator counts sources with a background, and the numerator counts positive-width response overlaps. Both scorers give the same exceptions. Intervals resample sources within each row and are marginal 95\% intervals. A single style edit supplies a point estimate only.}
\label{tab:v19-selectivity-category}
\begin{tabular}{@{}lrrrrrr@{}}
\toprule
Edit category & Overlap & Negative & FM mean & FM interval & MI mean & MI interval\\
\midrule
Random & 4/4 & 0 & $0.0820$ & $[0.0143, 0.1496]$ & $0.0557$ & $[0.0066, 0.1378]$ \\
Replace & 5/5 & 1 & $0.0248$ & $[-0.0322, 0.0901]$ & $0.0239$ & $[-0.0054, 0.0629]$ \\
Add & 5/5 & 0 & $0.0526$ & $[0.0239, 0.0812]$ & $0.0207$ & $[0.0136, 0.0265]$ \\
Delete & 5/5 & 0 & $0.0482$ & $[0.0100, 0.0977]$ & $0.0451$ & $[0.0108, 0.0884]$ \\
Content & 4/4 & 0 & $0.0505$ & $[0.0341, 0.0618]$ & $0.0451$ & $[0.0303, 0.0560]$ \\
Pose & 4/4 & 2 & $0.0051$ & $[-0.0316, 0.0419]$ & $-0.0013$ & $[-0.0300, 0.0246]$ \\
Color & 5/5 & 0 & $0.0655$ & $[0.0246, 0.1067]$ & $0.0483$ & $[0.0185, 0.0784]$ \\
Material & 4/5 & 0 & $0.0211$ & $[0.0105, 0.0316]$ & $0.0131$ & $[0.0022, 0.0240]$ \\
Background & 4/4 & 1 & $0.0006$ & $[-0.2419, 0.1707]$ & $0.0474$ & $[-0.0024, 0.0973]$ \\
Style & 1/1 & 0 & $0.0086$ & --- & $0.0043$ & --- \\
\midrule
All & 41/42 & 4 & $0.0390$ & $[0.0106, 0.0632]$ & $0.0326$ & $[0.0204, 0.0462]$ \\
\bottomrule
\end{tabular}
\end{table*}

The feature-map and masked-input scorers agree on the sign of all 41 supported
uniform--spatial contrasts (Table~\ref{tab:v19-selectivity-category}).
Two of four pose edits favor uniform: changing a dog's posture from lying
down to jumping ($-0.0119$ feature-map LPIPS) and opening a closed lid
($-0.0514$). The other exceptions replace a mountain reflection with a
ship reflection ($-0.0674$) and change a rice-field background ($-0.3656$).
The pose means are 0.0051 and $-0.0013$ under feature-map and masked-input
scoring, respectively. Background edits have means of 0.0006 and 0.0474.
Category support ranges from one to five images. These summaries locate
the observed reversals and quantify their sensitivity to the scorer.

Across all supported sources, the mean FlowEdit-minus-spatial contrast is
0.1196 [0.0739, 0.1631] under feature-map scoring and 0.0905 [0.0625, 0.1227]
under masked-input scoring. All 38 source signs agree. For FlowEdit minus
uniform, the corresponding estimates are 0.0638 [0.0229, 0.1046] and
0.0489 [0.0220, 0.0790]. The scorers agree on 39 of 41 source signs.
Figure~\ref{fig:v19-response-sensitivity} gives the same-source default and matched
gaps, the three paired distributions and the mass-weighted localization plane.
